\documentclass[11pt,a4paper]{article}

% ── Packages ──────────────────────────────────────────────────────────────────
\usepackage[margin=2.5cm]{geometry}
\usepackage{amsmath,amssymb}
\usepackage{booktabs}
\usepackage{hyperref}
\usepackage{xcolor}
\usepackage{longtable}
\usepackage{array}
\usepackage{graphicx}
\usepackage{fancyhdr}
\usepackage{titlesec}
\usepackage{enumitem}

% ── Colour definitions ────────────────────────────────────────────────────────
\definecolor{apexblue}{RGB}{0,51,102}
\definecolor{apexgray}{RGB}{100,100,100}
\definecolor{apexred}{RGB}{180,30,30}

% ── Section formatting ────────────────────────────────────────────────────────
\titleformat{\section}{\large\bfseries\color{apexblue}}{}{0em}{\thesection\quad}
\titleformat{\subsection}{\normalsize\bfseries\color{apexblue}}{}{0em}{\thesubsection\quad}

% ── Header / footer ───────────────────────────────────────────────────────────
\pagestyle{fancy}
\fancyhf{}
\rhead{\textcolor{apexgray}{\small RESTRICTED --- Model Risk Sensitive}}
\lhead{\textcolor{apexgray}{\small APEX-MDL-0012 \textbar{} MVA MDD v1.0}}
\rfoot{\textcolor{apexgray}{\small \thepage}}
\lfoot{\textcolor{apexgray}{\small Apex Global Bank \textbar{} Quantitative Research Division}}
\renewcommand{\headrulewidth}{0.4pt}
\renewcommand{\footrulewidth}{0.4pt}

% ── Hyperref setup ────────────────────────────────────────────────────────────
\hypersetup{
  colorlinks=true,
  linkcolor=apexblue,
  urlcolor=apexblue,
  pdftitle={Apex Global Bank -- MVA Model Development Document},
  pdfauthor={XVA Quant Team}
}

% ══════════════════════════════════════════════════════════════════════════════
\begin{document}
% ══════════════════════════════════════════════════════════════════════════════

\begin{titlepage}
  \begin{center}
    \vspace*{2cm}
    {\color{apexblue}\rule{\linewidth}{2pt}}\\[0.5cm]
    {\Huge\bfseries\color{apexblue} Model Development Document}\\[0.4cm]
    {\LARGE\bfseries Margin Valuation Adjustment (MVA)}\\[0.3cm]
    {\Large Version 1.0}\\[0.5cm]
    {\color{apexblue}\rule{\linewidth}{2pt}}\\[1cm]
    {\large Apex Global Bank \textbar{} Quantitative Research Division}\\[0.3cm]
    {\large\textcolor{apexred}{Classification: RESTRICTED --- Model Risk Sensitive}}\\[0.3cm]
    {\large SR 11-7 Section Reference: Model Development}\\[2cm]
    \begin{tabular}{ll}
      \textbf{Model ID:}   & APEX-MDL-0012 \\[4pt]
      \textbf{Date:}       & March 2026 \\[4pt]
      \textbf{Status:}     & In Validation \\[4pt]
      \textbf{Risk Tier:}  & Tier 1 \\
    \end{tabular}
  \end{center}
\end{titlepage}

% ── Table of Contents ─────────────────────────────────────────────────────────
\tableofcontents
\newpage

% ══════════════════════════════════════════════════════════════════════════════
\section{Document Metadata}
% ══════════════════════════════════════════════════════════════════════════════

\begin{table}[h!]
\centering
\caption{Document Control Information}
\begin{tabular}{@{}ll@{}}
\toprule
\textbf{Field} & \textbf{Value} \\
\midrule
Model Name            & Margin Valuation Adjustment (MVA) \\
Model ID              & APEX-MDL-0012 \\
Version               & 1.0 \\
Status                & In Validation \\
Risk Tier             & Tier 1 \\
Regulatory Framework  & BCBS-IOSCO UMR; ISDA SIMM v2.6; EMIR; Dodd-Frank \\
Date of Development   & March 2026 \\
Business Owner        & Head of Global Markets Trading (James Okafor) \\
Model Developer       & Quantitative Research --- XVA Team \\
MRO Review            & Pending --- Dr.\ Rebecca Chen \\
CRO Approval          & Pending --- Dr.\ Priya Nair \\
\bottomrule
\end{tabular}
\end{table}

\subsection{Version History}

\begin{table}[h!]
\centering
\caption{Document Version History}
\begin{tabular}{@{}lll@{}}
\toprule
\textbf{Version} & \textbf{Date} & \textbf{Change Summary} \\
\midrule
0.1 & Nov 2025 & Initial framework --- IM forecasting via regression \\
0.5 & Jan 2026 & SIMM sensitivities methodology integrated \\
1.0 & Mar 2026 & Production candidate; submitted for validation \\
\bottomrule
\end{tabular}
\end{table}

% ══════════════════════════════════════════════════════════════════════════════
\section{Executive Summary}
% ══════════════════════════════════════════════════════════════════════════════

\subsection{Model Purpose}

MVA is the present value of the expected future costs of posting Initial Margin (IM) on
uncleared OTC derivatives over the life of a trade. It was introduced as a necessary valuation
adjustment following the implementation of the BCBS-IOSCO Uncleared Margin Rules (UMR), which
began phased implementation in 2016 and reached Phase~6 in September 2022 (capturing entities
with AANA $\geq$ \$8 billion).

Unlike Variation Margin (which is exchanged daily and represents current mark-to-market),
Initial Margin is a forward-looking risk buffer held by a third-party custodian to cover
potential exposure during the Margin Period of Risk if a counterparty defaults. It is not
returned to the poster until the trade matures or is terminated. The cost of funding this
IM at Apex's unsecured borrowing spread for potentially decades is the MVA:

\begin{equation}
\mathrm{MVA} = \mathrm{PV}\!\left[\sum_i s_f(t_i)\cdot\mathbb{E}\!\left[\mathrm{IM}_{\mathrm{SIMM}}(t_i)\right]\cdot\Delta t_i\right]
\tag{1}
\end{equation}

\subsection{Why MVA Matters Now}

Before UMR Phase~6, only the largest dealers exchanged IM on bilateral OTC trades. Phase~6
extended the requirement to mid-size institutions. For Apex, this created an overnight
obligation to post IM on approximately 47\% of its bilateral OTC derivatives portfolio with
Phase~6 counterparties. The aggregate IM obligation for Apex is currently approximately
\textbf{\$8.7 billion} across all bilateral relationships.

At an average funding spread of 55\,bp, the annual cost of carrying this IM is approximately
\textbf{\$48 million per year}. Over the average remaining life of the portfolio (approximately
6.2 years duration-weighted), the present value of these future funding costs is approximately
\textbf{\$240 million} --- the current portfolio MVA.

\subsection{Output Description}

\begin{itemize}[leftmargin=*]
  \item \textbf{Trade MVA}: MVA charge per new trade --- incorporated into front-office pricing.
  \item \textbf{Portfolio MVA}: Daily P\&L reported in income statement.
  \item \textbf{IM forecast profile}: Expected IM at each future date (EPE of IM under SIMM).
  \item \textbf{MVA Greeks}: Sensitivities to rates and volatility for hedging.
  \item \textbf{MVA by counterparty}: Attribution for collateral efficiency analysis.
  \item \textbf{SIMM sensitivity report}: The underlying SIMM sensitivities driving IM calculations.
\end{itemize}

\subsection{Key Assumptions}

\begin{enumerate}[leftmargin=*]
  \item Initial Margin is calculated under ISDA SIMM v2.6 (the industry standard IM model for bilateral trades).
  \item Future IM is estimated by simulating future SIMM sensitivities (delta, vega, curvature) under the same Monte Carlo framework used for CVA/FVA.
  \item Apex's funding spread for IM is 55\,bp over OIS for 1--5Y; funding curve used for longer tenors.
  \item The custodian is a qualifying third-party custodian (State Street); IM posted earns no return (or OIS $-$ haircut).
  \item SIMM risk weights and correlations are fixed at current ISDA SIMM v2.6 calibration; no scenario for SIMM recalibration risk.
\end{enumerate}

\subsection{Known Limitations}

\begin{itemize}[leftmargin=*]
  \item Forecasting SIMM sensitivities under future market scenarios is computationally expensive and relies on linear approximation for efficiency; true future SIMM may differ by 5--15\%.
  \item SIMM recalibration risk: ISDA recalibrates SIMM risk weights annually; a recalibration that increases risk weights would increase IM requirements and thus MVA.
  \item The model assumes continuous IM posting; in practice, MTAs create discontinuities that are smoothed over in the simulation.
  \item MVA is computed assuming the full UMR regime; any regulatory change (e.g., threshold increases) would require immediate model update.
\end{itemize}

% ══════════════════════════════════════════════════════════════════════════════
\section{Business Context and Purpose}
% ══════════════════════════════════════════════════════════════════════════════

\subsection{The UMR Backstory}

The Uncleared Margin Rules emerged from the G20's Pittsburgh Summit (2009) mandate to reform
OTC derivatives markets following the 2008 crisis. The core insight: the daisy chain of
uncollateralised bilateral OTC trades had created a web of hidden counterparty risk. AIG's
near-failure was a direct result --- AIG had written enormous quantities of credit default swaps
with minimal collateral, and when the credit markets collapsed, the potential losses were so
large that AIG would have defaulted on hundreds of counterparties simultaneously.

UMR forces dealers and large end-users to post two-way IM on bilateral trades --- creating a
genuine buffer that survives the default of either party. The cost of this buffer --- MVA ---
is real and must be priced.

\subsection{ISDA SIMM Overview}

ISDA's Standard Initial Margin Model (SIMM) is the industry-standard methodology for
calculating IM on bilateral uncleared OTC derivatives. It was developed collaboratively by
major dealers and ISDA, approved by regulators, and adopted by essentially all Phase~1--6
participants.

SIMM computes IM as a risk-weighted sum of trade sensitivities. For a single risk class the
formula is:

\begin{equation}
\mathrm{IM}_{\mathrm{SIMM}} = \sqrt{\sum_i \sum_j \rho_{ij}\cdot \mathrm{WS}_i\cdot\mathrm{WS}_j}
\tag{2}
\end{equation}

where $\mathrm{WS}_i = \mathrm{RW}_i \cdot s_i$ (risk-weighted sensitivity), $\rho_{ij}$ is
the prescribed inter-bucket correlation, and $\mathrm{RW}_i$ is the SIMM risk weight for
bucket~$i$. SIMM covers six risk classes: Interest Rate (IR), Credit Qualifying (CQ), Credit
Non-Qualifying (CNQ), Equity (EQ), FX, and Commodity (CM).

\subsection{Why MVA Is Tier 1}

MVA feeds directly into:
\begin{enumerate}[leftmargin=*]
  \item \textbf{Income statement}: Daily MVA P\&L affects earnings.
  \item \textbf{New trade pricing}: Incorrect MVA pricing leads to systematic mis-pricing of bilateral derivatives.
  \item \textbf{Collateral optimisation decisions}: MVA drives decisions about which trades to clear (CCPs require IM too, but on different terms) vs.\ keep bilateral.
\end{enumerate}

% ══════════════════════════════════════════════════════════════════════════════
\section{Theoretical Foundation}
% ══════════════════════════════════════════════════════════════════════════════

\subsection{MVA Formula}

\begin{equation}
\mathrm{MVA} = -\sum_i s_f(t_i)\cdot\mathbb{E}\!\left[\mathrm{IM}_{\mathrm{SIMM}}(t_i)\right]\cdot D(0,t_i)\cdot\Delta t_i
\tag{3}
\end{equation}

where:
\begin{itemize}[leftmargin=*]
  \item $s_f(t_i)$ = Apex's funding spread at time $t_i$ (from the funding curve used in FVA)
  \item $\mathbb{E}[\mathrm{IM}_{\mathrm{SIMM}}(t_i)]$ = expected Initial Margin under SIMM at future time $t_i$
  \item $D(0,t_i)$ = OIS discount factor
  \item The negative sign: IM is a cash outflow (cost to Apex)
\end{itemize}

The key challenge is computing $\mathbb{E}[\mathrm{IM}_{\mathrm{SIMM}}(t_i)]$ --- forecasting
future SIMM sensitivities and therefore future IM requirements.

\subsection{Forecasting Future SIMM Sensitivities}

We employ three approaches of increasing accuracy and computational cost.

\textbf{Approach 1 --- Regression (current production approach)}: A regression model maps
current portfolio sensitivities and market state to future SIMM:

\begin{equation}
\mathrm{IM}_{\mathrm{SIMM}}(t) \approx f\!\left(\mathbf{S}(0),\, \mathbf{X}(t)\right)
\tag{4}
\end{equation}

where $\mathbf{S}(0)$ denotes current DV01/vega/curvature and $\mathbf{X}(t)$ the projected
market state. Model: gradient boosting regression trained on historical SIMM calculations.
Accuracy: within 8\% of full-calculation SIMM for 90\% of scenarios. Computation: $\approx$2
seconds per counterparty netting set.

\textbf{Approach 2 --- Approximate SIMM re-calculation}: At each Monte Carlo time step,
re-compute approximate SIMM using perturbed sensitivities via first-order Taylor expansion:

\begin{equation}
s_k(t) \approx s_k(0) + \frac{\partial s_k}{\partial r}\,\Delta r(t) + \frac{\partial s_k}{\partial\sigma}\,\Delta\sigma(t)
\tag{5}
\end{equation}

Accuracy: within 5\% of full SIMM for vanilla portfolios. Computation: $\approx$15 seconds per
counterparty.

\textbf{Approach 3 --- Full SIMM re-calculation (benchmark only)}: Re-run full SIMM at each
simulation node. Used only for quarterly benchmark validation.

\textbf{Current production}: Approach~1 for most counterparties; Approach~2 for the 20 largest
counterparties by IM obligation (representing $\approx$75\% of total MVA).

\subsection{SIMM Sensitivities Under Future Market Scenarios}

For the interest rate risk class (62\% of total IM at Apex), delta sensitivity evolution
assumes future DV01 scales with remaining duration and the simulated rate environment under a
two-factor Hull-White model:

\begin{equation}
\mathrm{DV01}(t) \approx \mathrm{DV01}(0)\cdot\frac{D_{\mathrm{rem}}(t)}{D(0)}\cdot\Phi_r(t)
\tag{6}
\end{equation}

where $\Phi_r(t)$ is a rate-level adjustment factor derived from the simulated short rate.
Vega (for swaption portfolios) is proportional to remaining time to expiry and the ratio of
current vs.\ future implied volatility. Vega collapses to zero at option expiry.

\subsection{Marginal MVA for New Trade Pricing}

When pricing a new trade, the relevant quantity is the \textbf{marginal MVA} --- the change in
portfolio MVA from adding the new trade:

\begin{equation}
\mathrm{MVA}_{\mathrm{marginal}} = \mathrm{MVA}(\text{Portfolio} + \text{New Trade}) - \mathrm{MVA}(\text{Portfolio})
\tag{7}
\end{equation}

Due to netting effects within a counterparty's netting set, marginal MVA can be significantly
less than the standalone MVA of the new trade, or even negative (IM-reducing) if the new trade
offsets existing sensitivities.

\subsection{Literature References}

\begin{itemize}[leftmargin=*]
  \item ISDA (2022). \textit{SIMM Methodology, version 2.6}. ISDA.
  \item BCBS-IOSCO (2015). \textit{Margin Requirements for Non-Centrally Cleared Derivatives}, rev.\ 2020.
  \item Green, A., Kenyon, C.\ (2014). ``MVA: Initial Margin Valuation Adjustment by Replication and Regression.'' \textit{Risk Magazine}, May 2014.
  \item Anfuso, F., Aziz, D., Giltinan, P., Loukopoulos, K.\ (2017). ``A Sound Modelling and Backtesting Framework for Forecasting Initial Margin Requirements.'' SSRN Working Paper.
  \item Caspers, P., Giltinan, P., Lichters, R., Nowaczyk, N.\ (2017). ``Forecasting Initial Margin Requirements: A Model Evaluation.'' \textit{Journal of Risk}, Vol.\ 19.
\end{itemize}

% ══════════════════════════════════════════════════════════════════════════════
\section{Data Requirements and Governance}
% ══════════════════════════════════════════════════════════════════════════════

\subsection{Input Data Sources}

\begin{table}[h!]
\centering
\caption{MVA Model Input Data Sources}
\begin{tabular}{@{}lll@{}}
\toprule
\textbf{Data Item} & \textbf{Source} & \textbf{Frequency} \\
\midrule
SIMM risk weights and correlations & ISDA SIMM documentation (annual) & Annual; immediate on release \\
Trade sensitivities (delta, vega, curvature) & Risk systems (Murex) & Daily \\
CSA terms (IM threshold, MTA, eligible collateral) & ISDA SIMM Custodial Agreement DB & On amendment \\
Funding curve & Treasury (same as FVA) & Daily \\
Historical SIMM calculations (for regression) & Risk data warehouse & Monthly re-training \\
Actual IM calls received/posted & Collateral management system & Daily \\
\bottomrule
\end{tabular}
\end{table}

\subsection{SIMM Version Governance}

ISDA releases updated SIMM calibrations annually (typically September). Apex's process:
\begin{enumerate}[leftmargin=*]
  \item ISDA releases SIMM vX.Y (typically 60 days notice).
  \item XVA team recalibrates MVA model to new SIMM weights.
  \item Model Risk Officer expedited review (15 business days for SIMM version updates).
  \item Production deployment by SIMM effective date.
  \item Parallel run for 5 business days pre-cutover.
\end{enumerate}

\textbf{SIMM version currently in production}: v2.6 (effective September 2024).

\subsection{BCBS 239 Data Lineage}

Every MVA figure is traceable to: (a)~trade data from Murex, (b)~SIMM sensitivity
calculation (timestamped), (c)~SIMM version tag, (d)~funding curve snapshot, (e)~MVA
model version.

% ══════════════════════════════════════════════════════════════════════════════
\section{Methodology and Implementation}
% ══════════════════════════════════════════════════════════════════════════════

\subsection{Production Run Architecture}

MVA shares the simulation infrastructure with CVA and FVA. The additional computation is:
\begin{enumerate}[leftmargin=*]
  \item At each Monte Carlo time step, compute approximate future SIMM for each netting set.
  \item Multiply by funding spread at that tenor.
  \item Discount back to today.
\end{enumerate}

\textbf{Run frequency}: Daily EOD (10,000 paths, same as CVA). Intraday indicative runs every
4 hours (1,000 paths).

\textbf{Computation time}: Approximately 18 minutes incremental on top of CVA run (regression
approach). Total XVA suite (CVA $+$ FVA $+$ MVA) target: $\leq$ 70 minutes.

\subsection{IM Portfolio Optimisation}

MVA is used in the collateral optimisation engine to evaluate whether bilateral trades should be:
\begin{enumerate}[leftmargin=*]
  \item Kept bilateral (pay MVA and FVA; retain flexibility).
  \item Cleared at CCP (pay CCP IM and default fund contribution; standardised but efficient for plain vanilla).
  \item Compressed/terminated (reduce gross notional; reduce IM obligation).
\end{enumerate}

The optimisation runs weekly and generates recommendations for the Head of Operations and
trading desk heads.

\subsection{Actual IM vs.\ Model Validation}

Daily comparison of model-predicted SIMM against actual IM calls received from counterparties:
\begin{itemize}[leftmargin=*]
  \item Tolerance: $\leq 5\%$ difference for 95\% of daily IM calls.
  \item Systematic bias (model consistently under- or over-predicting): investigate within 3 business days.
\end{itemize}

This is both a data quality check (are counterparties calculating SIMM correctly?) and a
model accuracy check.

% ══════════════════════════════════════════════════════════════════════════════
\section{Model Testing and Backtesting}
% ══════════════════════════════════════════════════════════════════════════════

\subsection{Actual vs.\ Predicted IM}

Primary backtesting: compare model-generated expected IM at horizon $T$ against actual IM
posted at time $T$.

Sample size: 24 months of daily IM calls across 130+ active counterparties.

\begin{table}[h!]
\centering
\caption{Backtesting Results (as of Model Submission)}
\begin{tabular}{@{}ll@{}}
\toprule
\textbf{Metric} & \textbf{Result} \\
\midrule
Mean absolute error            & 4.2\% of actual IM (within 5\% threshold) \\
95th percentile error          & 11.8\% (outliers from SIMM recalibration events) \\
Directional accuracy           & 87\% \\
\bottomrule
\end{tabular}
\end{table}

\subsection{Full-Calculation Benchmark}

Quarterly comparison of regression-based SIMM forecast (production) against Approach~3 (full
SIMM re-calculation) for a representative sub-portfolio:
\begin{itemize}[leftmargin=*]
  \item Target: regression within 8\% of full calculation.
  \item Last benchmark result (Q4 2025): 6.3\% mean deviation --- \textbf{pass}.
\end{itemize}

\subsection{SIMM Sensitivity Stress Test}

Simulate a 50\% increase in all SIMM risk weights (analogous to a severe SIMM recalibration):

\begin{table}[h!]
\centering
\caption{SIMM Stress Test Results}
\begin{tabular}{@{}ll@{}}
\toprule
\textbf{Impact Item} & \textbf{Result} \\
\midrule
Portfolio IM increase         & \$4.3B (from \$8.7B to $\approx$\$13B) \\
MVA increase                  & $\approx$\$120M \\
Funding requirement increase  & $\approx$\$4.3B \\
LCR impact                    & Remains above 100\% minimum \\
\bottomrule
\end{tabular}
\end{table}

% ══════════════════════════════════════════════════════════════════════════════
\section{Performance Metrics and Calibration Standards}
% ══════════════════════════════════════════════════════════════════════════════

\begin{table}[h!]
\centering
\caption{MVA Performance Metrics and Thresholds}
\begin{tabular}{@{}lll@{}}
\toprule
\textbf{Metric} & \textbf{Acceptable Range} & \textbf{Action if Breached} \\
\midrule
IM prediction error vs.\ actual calls & $<5\%$ mean absolute & Investigate within 3 days \\
IM prediction error vs.\ full SIMM    & $<8\%$ mean deviation & Review regression model \\
SIMM version deployment lag           & $\leq 0$ days vs.\ ISDA date & Emergency deployment process \\
Actual IM call reconciliation         & $<3$ unreconciled calls $>$\$10M & Collateral ops escalation \\
MVA P\&L explain                      & $>75\%$ & Investigate within 2 days \\
\bottomrule
\end{tabular}
\end{table}

% ══════════════════════════════════════════════════════════════════════════════
\section{Limitations, Assumptions, and Known Weaknesses}
% ══════════════════════════════════════════════════════════════════════════════

\subsection{SIMM Recalibration Risk}

ISDA recalibrates SIMM annually. A significant recalibration would:
\begin{enumerate}[leftmargin=*]
  \item Increase actual IM requirements immediately.
  \item Require MVA model update within 60 days.
  \item Affect profitability of all existing bilateral derivatives --- a retroactive cost that cannot be hedged.
\end{enumerate}

No model-based compensation is possible for this risk. The compensating control is operational:
ISDA recalibration schedule is tracked; pre-implementation impact analysis of proposed changes
is prepared quarterly.

\subsection{Basis Between SIMM and Actual IM}

Different counterparties calculate SIMM slightly differently (different curve construction,
different aggregation methods). These differences create a basis between Apex's SIMM calculation
and the counterparty's --- resulting in disputes and IM calls that differ from model predictions.

\textbf{Compensating control}: Dispute resolution process (ISDA SIMM reconciliation guidance);
daily IM call reconciliation with automatic escalation for $>$\$5M discrepancy.

\subsection{Regression Model Staleness}

The regression model for future SIMM is retrained monthly on historical data. In a regime change
(e.g., a rapid large move in rates or a significant restructuring of the derivatives portfolio),
the regression may become unrepresentative until the next retraining cycle.

\textbf{Compensating control}: Anomaly detection on regression outputs --- if predicted IM
deviates from actual by $>15\%$ on three consecutive days, emergency retraining is triggered.

\subsection{Eligible Collateral Uncertainty}

SIMM IM must be posted in eligible collateral (typically cash or high-quality government bonds).
If eligible collateral is scarce (a real stress scenario in which everyone is simultaneously
posting IM), the cost of sourcing eligible collateral may exceed the funding spread assumed in
the model.

\textbf{Compensating control}: Collateral availability stress test run quarterly; minimum
eligible collateral buffer maintained by Treasury.

% ══════════════════════════════════════════════════════════════════════════════
\section{Compensating Controls}
% ══════════════════════════════════════════════════════════════════════════════

\begin{table}[h!]
\centering
\caption{MVA Compensating Controls}
\begin{tabular}{@{}p{4.5cm}p{6cm}p{3.5cm}@{}}
\toprule
\textbf{Risk} & \textbf{Control} & \textbf{Owner} \\
\midrule
SIMM recalibration & Pre-implementation impact analysis; operational monitoring & XVA Quant + Reg Affairs \\[4pt]
IM calculation disputes & Daily IM reconciliation; dispute resolution protocol & Collateral Ops \\[4pt]
Regression staleness & Monthly retraining; anomaly detection & XVA Quant Team \\[4pt]
Eligible collateral scarcity & Collateral availability stress test; buffer maintained & Treasury \\[4pt]
SIMM version deployment lag & Change management process; 60-day advance preparation & Technology / XVA \\
\bottomrule
\end{tabular}
\end{table}

\subsection{Model Reserve}

A model reserve of \textbf{\$18 million} is held against MVA model uncertainty:

\begin{table}[h!]
\centering
\caption{MVA Model Reserve Breakdown}
\begin{tabular}{@{}lr@{}}
\toprule
\textbf{Reserve Component} & \textbf{Amount (USD M)} \\
\midrule
Regression approximation error  & 8 \\
SIMM recalibration risk         & 7 \\
Collateral cost uncertainty     & 3 \\
\midrule
\textbf{Total Reserve}          & \textbf{18} \\
\bottomrule
\end{tabular}
\end{table}

% ══════════════════════════════════════════════════════════════════════════════
\section{Change Management}
% ══════════════════════════════════════════════════════════════════════════════

\subsection{SIMM Version Updates}

SIMM annual recalibration is classified as a material model change requiring expedited MRO
review (15 business days). No delay to production deployment is permitted --- ISDA effective
dates are regulatory compliance deadlines.

\subsection{All Other Material Changes}

Full re-validation process applies (definition same as CVA MDD Section~11). Material changes
include: change to IM forecasting approach (Approach~1/2/3 transition), change to the
shared simulation infrastructure, change to the funding curve used in MVA computation,
change to eligible collateral assumptions.

% ══════════════════════════════════════════════════════════════════════════════
\section{Approvals and Sign-Off}
% ══════════════════════════════════════════════════════════════════════════════

\begin{table}[h!]
\centering
\caption{Approval Status}
\begin{tabular}{@{}llll@{}}
\toprule
\textbf{Role} & \textbf{Name} & \textbf{Status} & \textbf{Date} \\
\midrule
Model Developer                & XVA Quant Team    & Submitted & March 2026 \\
Model Risk Officer             & Dr.\ Rebecca Chen & Pending   & --- \\
Chief Risk Officer             & Dr.\ Priya Nair   & Pending   & --- \\
Head of Treasury (UMR Compliance) & Amara Diallo   & Pending   & --- \\
Business Owner                 & James Okafor      & Pending   & --- \\
\bottomrule
\end{tabular}
\end{table}

\vspace{1cm}
\noindent\textcolor{apexgray}{\footnotesize
Document Classification: RESTRICTED --- Model Risk Sensitive\\
Apex Global Bank \textbar{} Quantitative Research Division \textbar{} XVA Team\\
Next scheduled review: March 2027 (or upon SIMM version update / material model change)
}

\end{document}
